Los experimentos y mediciones de tiempo de esta Tarea 1 fueron ejecutados bajo las siguientes especificaciones técnicas:

\begin{itemize}
    \item \textbf{Hardware:} Procesador 13th Gen Intel(R) Core(TM) i5-1335U @ 1.30 GHz con 16 GB de memoria RAM.
    
    \item \textbf{Sistema Operativo:} Windows 11 Pro (versión 25H2) utilizando el subsistema Windows para Linux (WSL) con la distribución Ubuntu 24.04.2 LTS (con aproximadamente 7.6 GiB de RAM asignada) como entorno de ejecución principal.
    
    \item \textbf{Compilador y Herramientas:} Los algoritmos fueron implementados en C++ y compilados con G++ versión 13.3.0 dentro del entorno Linux. La generación de los casos de prueba y los gráficos de resultados se realizaron utilizando Python 3.12.3.
\end{itemize}

Para la evaluación empírica, se utilizaron los datasets sintéticos descritos en el enunciado, generados mediante los scripts proporcionados en la carpeta de la tarea. Para los \textbf{algoritmos de ordenamiento} se generaron arreglos de tamaño $n$ para analizar el comportamiento temporal de los algoritmos (Mergesort, Patiencesort, Quicksort y Sort nativo de C++) bajo tres condiciones iniciales: datos generados aleatoriamente, ordenados de forma ascendente y de forma descendente. Por otro lado, para la \textbf{multiplicación de matrices}, se utilizaron matrices cuadradas de $n \times n$ pobladas con números enteros aleatorios para comparar el rendimiento temporal y espacial del algoritmo clásico (Naive) frente al algoritmo de Strassen.